\chapter{Thực nghiệm và đánh giá kết quả}
\label{chap:experiments}
%=================================================

\section{Kiến trúc hệ thống HybridFallTransformer}
\label{sec:proposed_system}

Phần này trình bày chi tiết kiến trúc hệ thống HybridFallTransformer -- hệ thống phát hiện té ngã end-to-end được đề xuất trong khóa luận này. Hệ thống được thiết kế theo nguyên tắc privacy-by-design, hoạt động thời gian thực trên thiết bị biên và đạt độ chính xác cao trong việc phân biệt té ngã với các hoạt động sinh hoạt thông thường.

\subsection{Pipeline phát hiện té ngã}
\label{sec:pipeline_overview}

Pipeline tổng quan của hệ thống HybridFallTransformer được minh họa trong Hình \ref{fig:pipeline_overview}. Luồng xử lý bao gồm năm giai đoạn chính:

\begin{enumerate}
    \item \textbf{Thu nhận video} từ camera giám sát RGB.
    \item \textbf{Trích xuất khung xương} bằng mô hình YOLOv11-Pose, cho ra vector tọa độ 17 điểm khớp COCO.
    \item \textbf{Biến đổi đặc trưng} bằng module PIFR, tạo vector 60 chiều cho mỗi khung hình.
    \item \textbf{Mô hình hóa chuỗi thời gian} bằng Hybrid Transformer Encoder với Temporal Self-Attention.
    \item \textbf{Phân loại và cảnh báo} -- nếu phát hiện té ngã, hệ thống gửi thông báo qua Telegram.
\end{enumerate}

%\begin{figure}[h]
%    \centering
%    \includegraphics[width=0.95\textwidth]{figures/chapter3/pipeline_overview.png}
%    \caption{Sơ đồ pipeline hệ thống HybridFallTransformer}
%    \label{fig:pipeline_overview}
%\end{figure}

Điểm nổi bật của pipeline này là toàn bộ dữ liệu hình ảnh RGB chỉ được sử dụng tại bước 2 để trích xuất khung xương, sau đó bị hủy ngay lập tức. Các bước 3-5 chỉ làm việc với vector tọa độ số, đảm bảo quyền riêng tư tuyệt đối cho đối tượng được giám sát.

\subsection{Module trích xuất khung xương YOLOv11-Pose}
\label{sec:yolo_pose_module}

YOLOv11-Pose là biến thể của YOLOv11 được tối ưu cho bài toán ước lượng tư thế người. Module này nhận đầu vào là khung hình RGB và đầu ra là danh sách các người được phát hiện cùng với tọa độ 17 điểm khớp COCO cho mỗi người.

Cấu hình của module trong hệ thống:

\begin{itemize}
    \item \textbf{Kích thước đầu vào}: 640 × 640 pixel
    \item \textbf{Ngưỡng confidence}: 0.5
    \item \textbf{Ngưỡng NMS}: 0.5
    \item \textbf{Kết quả đầu ra}: tensor có shape $(N, 17, 3)$ với $N$ là số người trong khung hình, mỗi điểm gồm $(x, y, confidence)$
\end{itemize}

Trong triển khai thực tế cho bài toán phát hiện té ngã, hệ thống giả định chỉ có một người trong khung hình (phù hợp với phòng bệnh nhân hoặc phòng ngủ người cao tuổi). Nếu phát hiện nhiều hơn một người, hệ thống chọn người có tổng confidence cao nhất.

\subsection{Module biến đổi đặc trưng PIFR}
\label{sec:pifr_module_detail}

Module PIFR (\textit{Pose-based Interesting Feature Representation}) được áp dụng cho mỗi khung hình sau khi trích xuất khung xương. Đầu vào của module là vector tọa độ 17 điểm × 3 giá trị $(x, y, conf)$, đầu ra là vector đặc trưng 60 chiều.

Quy trình biến đổi bao gồm bốn nhóm đặc trưng:

\begin{enumerate}
    \item \textbf{Đặc trưng vị trí tuyệt đối (10 chiều)}: tọa độ chuẩn hóa của đầu, vai, hông, trung tâm khung xương.
    
    \item \textbf{Đặc trưng khoảng cách xương (16 chiều)}: độ dài của các xương chính như vai, hông, tay, chân.
    
    \item \textbf{Đặc trưng góc (12 chiều)}: góc tại các khớp then chốt -- vai trái, vai phải, hông trái, hông phải, gáy, thân. Các góc này đặc biệt nhạy với trạng thái té ngã vì khi ngã, thân người thường tạo góc bất thường với mặt phẳng đứng.
    
    \item \textbf{Đặc trưng động học (22 chiều)}: tốc độ thay đổi (velocity) và gia tốc (acceleration) của 4 góc và 4 khoảng cách quan trọng nhất. Nhóm đặc trưng này phản ánh động lực học chuyển động -- té ngã có gia tốc lớn và đột ngột ở giai đoạn va chạm.
\end{enumerate}

Module PIFR không chứa tham số học được, hoàn toàn dựa trên tính toán hình học. Thời gian xử lý dưới 1ms trên CPU.

\subsection{Bộ mã hóa Hybrid Transformer}
\label{sec:hybrid_transformer_detail}

Sau khi biến đổi bằng PIFR, chuỗi vector đặc trưng 60D được đưa vào bộ mã hóa Hybrid Transformer để mô hình hóa phụ thuộc thời gian. Kiến trúc chi tiết được minh họa trong Hình \ref{fig:transformer_architecture}.

Kiến trúc Hybrid Transformer Encoder bao gồm các thành phần:

\textbf{Lớp nhúng đặc trưng (Feature Embedding Layer).} Vector 60D từ PIFR được ánh xạ lên không gian $d_{\text{model}}$ chiều thông qua một lớp Linear với LayerNorm:

\begin{equation}
z_0 = \text{LayerNorm}(W_{\text{embed}} \cdot f + b_{\text{embed}})
\end{equation}

trong đó $f \in \mathbb{R}^{60}$ là vector đặc trưng từ PIFR.

\textbf{Position Encoding.} Do Self-Attention là phép toán permutation-invariant, hệ thống cộng position encoding vào mỗi khung hình để mã hóa thông tin thứ tự thời gian. Sử dụng Sinusoidal Position Encoding.

\textbf{Temporal Self-Attention.} Mỗi Transformer layer chứa một Multi-Head Self-Attention:

\begin{equation}
\text{MultiHead}(Q, K, V) = \text{Concat}(\text{head}_1, \ldots, \text{head}_h) W^O
\end{equation}

với mỗi head tính:
\begin{equation}
\text{head}_h = \text{Attention}(Q W_q^h, K W_k^h, V W_v^h) = \text{softmax}\left(\frac{Q K^T}{\sqrt{d_k}}\right) V
\end{equation}

\textbf{Feed-Forward Network.} Sau attention, một mạng truyền thẳng hai lớp áp dụng:

\begin{equation}
\text{FFN}(x) = W_2 \cdot \text{GELU}(W_1 \cdot x + b_1) + b_2
\end{equation}

với GELU là activation function: $\text{GELU}(x) = x \cdot \Phi(x)$ (với $\Phi$ là CDF phân bố chuẩn).

\textbf{Transformer Layer.} Mỗi Transformer layer có cấu trúc:
\begin{align}
z' &= \text{LayerNorm}(z + \text{MultiHead}(z, z, z)) \\
z'' &= \text{LayerNorm}(z' + \text{FFN}(z'))
\end{align}

\textbf{Toán tử Mean Pooling.} Thay vì dùng token [CLS] như trong BERT, hệ thống sử dụng Mean Pooling để tổng hợp biểu diễn từ tất cả các khung hình:

\begin{equation}
\mathbf{z}_{\text{pooled}} = \frac{1}{T} \sum_{t=1}^{T} z_t^{(L)}
\end{equation}

\textbf{Lớp phân loại.} Biểu diễn sau Mean Pooling được đưa qua lớp phân loại:

\begin{align}
\hat{y} &= \sigma(W_{\text{cls}} \cdot \mathbf{z}_{\text{pooled}} + b_{\text{cls}})
\end{align}

với $\sigma$ là sigmoid activation cho phân loại nhị phân.

Tổng hợp số lượng tham số của mô hình:

\begin{table}[h]
\centering
\caption{Số lượng tham số của Hybrid Transformer}
\label{tab:transformer_params}
\begin{tabular}{|l|c|}
\hline
\textbf{Thành phần} & \textbf{Số tham số} \\
\hline
Feature Embedding & tùy d\_model \\
\hline
Position Encoding & 0 (fixed) \\
\hline
Multi-Head Attention (×N layers) & tùy nhead, d\_model \\
\hline
Feed-Forward Network (×N layers) & tùy d\_model \\
\hline
LayerNorm & tùy d\_model \\
\hline
Classifier & 1 \\
\hline
\end{tabular}
\end{table}

Với cấu hình sử dụng trong thực nghiệm ($d_{model}=256, nhead=4, num\_layers=3$), tổng số tham số của Hybrid Transformer là khoảng 1.8M, phù hợp để huấn luyện trên GPU.

%-------------------------------------------------

\section{Thiết lập thực nghiệm}
\label{sec:experimental_setup}

\subsection{Môi trường phát triển và phần cứng}
\label{sec:dev_environment}

Thực nghiệm được tiến hành trên hai môi trường:

\begin{itemize}
    \item \textbf{Máy tính phát triển}: Kaggle Notebook (GPU P100)
    \item \textbf Thiết bị biên}: Raspberry Pi 4 Model B (4GB)
\end{itemize}

Các thư viện và framework sử dụng:

\begin{itemize}
    \item Python 3.10
    \item PyTorch 2.1.0
    \item Ultralytics (YOLOv11) 8.0.200
    \item OpenCV 4.8.0
    \item NumPy 1.24.0
    \item scikit-learn 1.3.0
\end{itemize}

\subsection{Bộ dữ liệu}
\label{sec:dataset}

Thực nghiệm sử dụng kết hợp hai bộ dữ liệu chuẩn trong lĩnh vực phát hiện té ngã:

\begin{itemize}
    \item \textbf{CaucaFall Dataset} \cite{caucafall}: 100 video từ camera góc nhìn cao, bao gồm các sự kiện té ngã và hoạt động sinh hoạt thông thường.
    \item \textbf{MCFD (Multiple Cameras Fall Dataset)} \cite{mcfd}: 551 video từ nhiều góc camera khác nhau, bao gồm 04 loại té ngã và 08 loại hoạt động sinh hoạt thông thường (ADL).
\end{itemize}

Tổng hợp hai bộ dữ liệu, tập dữ liệu cuối cùng bao gồm:

\begin{itemize}
    \item \textbf{Tổng số mẫu}: 651 chuỗi khung xương
    \item \textbf{Số lượng Fall}: 313 mẫu (48.1\%)
    \item \textbf{Số lượng No Fall}: 338 mẫu (51.9\%)
    \item \textbf{Kích thước đặc trưng}: X=(651, 60, 60) -- 60 frames × 60 features
\end{itemize}

\begin{itemize}
    \item \textbf{Training set (80\%)}: 521 mẫu
    \item \textbf{Validation set (20\% của training)}: 52 mẫu
    \item \textbf{Test set (20\% mỗi fold)}: 130 mẫu mỗi fold trong 5-Fold CV
\end{itemize}

\begin{figure}[h]
    \centering
    \begin{minipage}{0.45\textwidth}
        \centering
        %\includegraphics[width=\textwidth]{figures/chapter3/caucafall_sample.png}
        \caption[Mẫu từ CaucaFall Dataset]{Mẫu từ CaucaFall Dataset}
        \label{fig:caucafall_sample}
    \end{minipage}
    \hfill
    \begin{minipage}{0.45\textwidth}
        \centering
        %\includegraphics[width=\textwidth]{figures/chapter3/mcfd_sample.png}
        \caption[Mẫu từ MCFD Dataset]{Mẫu từ MCFD Dataset}
        \label{fig:mcfd_sample}
    \end{minipage}
\end{figure}

\subsection{Cấu hình huấn luyện}
\label{sec:training_config}

Bảng \ref{tab:training_config} tổng hợp các siêu tham số huấn luyện:

\begin{table}[h]
\centering
\caption{Cấu hình huấn luyện HybridFallTransformer}
\label{tab:training_config}
\begin{tabular}{|l|c|}
\hline
\textbf{Hyperparameter} & \textbf{Giá trị} \\
\hline
Optimizer & AdamW \\
\hline
Learning rate & 5e-4 \\
\hline
Weight decay & 1e-4 \\
\hline
Batch size & 32 \\
\hline
Epochs & 100 \\
\hline
Scheduler & Cosine Annealing \\
\hline
Early stopping patience & 15 epochs \\
\hline
Dropout & 0.2 \\
\hline
$d_{model}$ & 256 \\
\hline
Number of heads & 4 \\
\hline
Number of layers & 3 \\
\hline
Number of folds & 5 \\
\hline
\end{tabular}
\end{table}

Mô hình được huấn luyện với batch size 32 trong tối đa 100 epochs. Early stopping được áp dụng với patience 15 epochs -- nếu validation loss không cải thiện trong 15 epochs liên tiếp, quá trình huấn luyện dừng lại và khôi phục trọng số tốt nhất.

Hàm mất mát sử dụng là Binary Cross-Entropy with Logits:

\begin{equation}
\mathcal{L} = -\frac{1}{N} \sum_{i=1}^{N} \left[ y_i \log(\sigma(\hat{y}_i)) + (1 - y_i) \log(1 - \sigma(\hat{y}_i)) \right]
\end{equation}

\subsection{Các độ đo đánh giá}
\label{sec:metrics}

Hệ thống được đánh giá dựa trên các độ đo phân loại chuẩn:

\begin{align}
\text{Accuracy} &= \frac{TP + TN}{TP + TN + FP + FN} \\
\text{Precision} &= \frac{TP}{TP + FP} \\
\text{Recall (Sensitivity)} &= \frac{TP}{TP + FN} \\
\text{F1-score} &= \frac{2 \times \text{Precision} \times \text{Recall}}{\text{Precision} + \text{Recall}} \\
\text{Specificity} &= \frac{TN}{TN + FP}
\end{align}

Trong bài toán phát hiện té ngã:
\begin{itemize}
    \item \textbf{True Positive (TP)}: Mẫu Fall được phân loại đúng là Fall
    \item \textbf{True Negative (TN)}: Mẫu ADL được phân loại đúng là ADL
    \item \textbf{False Positive (FP)}: Mẫu ADL bị nhầm thành Fall (báo động sai)
    \item \textbf{False Negative (FN)}: Mẫu Fall bị bỏ sót thành ADL (bỏ sót)
\end{itemize}

Đặc biệt, hai độ đo then chốt cho hệ thống cảnh báo là Recall (tỷ lệ phát hiện đúng) và Specificity (tỷ lệ không báo động sai trên ADL).

%-------------------------------------------------

\section{Kết quả thực nghiệm}
\label{sec:results}

\subsection{Kết quả huấn luyện 5-Fold Cross-Validation}
\label{sec:training_results}

Quá trình huấn luyện sử dụng 5-Fold Cross-Validation với StratifiedKFold. Mỗi fold được huấn luyện độc lập với early stopping. Bảng \ref{tab:kfold_results} trình bày kết quả chi tiết từng fold:

\begin{table}[h]
\centering
\caption{Kết quả 5-Fold Cross-Validation}
\label{tab:kfold_results}
\begin{tabular}{|c|c|c|c|c|c|}
\hline
\textbf{Fold} & \textbf{Train} & \textbf{Val} & \textbf{Test} & \textbf{Acc} & \textbf{F1} \\
\hline
1 & 468 & 52 & 131 & 0.8321 & 0.8281 \\
\hline
2 & 469 & 52 & 130 & 0.8308 & 0.8167 \\
\hline
3 & 469 & 52 & 130 & 0.7769 & 0.7680 \\
\hline
4 & 469 & 52 & 130 & 0.8077 & 0.7899 \\
\hline
5 & 469 & 52 & 130 & 0.8154 & 0.8000 \\
\hline
\end{tabular}
\end{table}

Quá trình huấn luyện hội tụ sau khoảng 60-98 epochs tùy fold, với early stopping kích hoạt khi validation loss không cải thiện trong 15 epochs liên tiếp.

\begin{itemize}
    \item \textbf{Fold 1}: Early stopping at epoch 98, Val F1 đạt 0.8571
    \item \textbf{Fold 2}: Early stopping at epoch 78, Val F1 đạt 0.8302
    \item \textbf{Fold 3}: Early stopping at epoch 98, Val F1 đạt 0.8182
    \item \textbf{Fold 4}: Early stopping at epoch 72, Val F1 đạt 0.8889
    \item \textbf{Fold 5}: Early stopping at epoch 63, Val F1 đạt 0.8772
    \item \textbf{Best Fold}: Fold 1 với F1 = 0.8281
\end{itemize}

\begin{figure}[h]
    \centering
    \begin{minipage}{0.45\textwidth}
        \centering
        %\includegraphics[width=\textwidth]{figures/chapter3/training_loss.png}
        \caption[Đồ thị Loss qua các epoch]{Đồ thị Loss qua các epoch}
        \label{fig:training_loss}
    \end{minipage}
    \hfill
    \begin{minipage}{0.45\textwidth}
        \centering
        %\includegraphics[width=\textwidth]{figures/chapter3/training_f1.png}
        \caption[Đồ thị F1-score trên validation]{Đồ thị F1-score trên validation}
        \label{fig:training_f1}
    \end{minipage}
\end{figure}

Hình \ref{fig:training_loss} cho thấy training loss giảm ổn định qua các epoch. Validation loss duy trì cùng xu hướng, không có dấu hiệu overfitting nghiêm trọng nhờ dropout và early stopping.

\subsection{Kết quả đánh giá tổng hợp 5-Fold CV}
\label{sec:aggregated_results}

Kết quả đánh giá được tổng hợp từ tất cả 5 folds, với confusion matrix được tính trên toàn bộ 651 predictions (mỗi sample được dự đoán đúng một lần trong fold tương ứng).

\begin{table}[h]
\centering
\caption{Kết quả tổng hợp 5-Fold Cross-Validation}
\label{tab:aggregated_results}
\begin{tabular}{|l|c|c|}
\hline
\textbf{Độ đo} & \textbf{Giá trị} & \textbf{Ghi chú} \\
\hline
Average Accuracy & 81.26\% & $\pm$ 2.21\% \\
\hline
Average F1-Score & 80.05\% & $\pm$ 2.22\% \\
\hline
Average AUC-ROC & 84.67\% & trên tất cả folds \\
\hline
\end{tabular}
\end{table}

\begin{figure}[hbtp]
    \centering
    %\includegraphics[width=0.7\textwidth]{figures/chapter3/confusion_matrix_5fold_cv.png}
    \caption[Confusion Matrix 5-Fold CV]{Confusion Matrix tổng hợp từ 5-Fold Cross-Validation}
    \label{fig:confusion_matrix_5fold}
\end{figure}

Hình \ref{fig:confusion_matrix_5fold} thể hiện Confusion Matrix tổng hợp từ 5-Fold Cross-Validation với 651 mẫu kiểm tra. Phân tích chi tiết:

\begin{itemize}
    \item \textbf{True Negative (TN = 292)}: 292 mẫu No Fall được phân loại đúng -- chiếm tỷ lệ cao trong tổng số mẫu No Fall.
    \item \textbf{False Positive (FP = 46)}: 46 mẫu No Fall bị nhầm thành Fall -- đây là các trường hợp báo động sai (false alarm).
    \item \textbf{False Negative (FN = 23)}: 23 mẫu Fall bị bỏ sót thành No Fall -- đây là các trường hợp nguy hiểm nhất vì té ngã không được phát hiện.
    \item \textbf{True Positive (TP = 290)}: 290 mẫu Fall được phát hiện đúng -- cho thấy mô hình có khả năng nhận diện té ngã khá tốt.
\end{itemize}

Từ Confusion Matrix, ta tính được:
\begin{align}
\text{Accuracy} &= \frac{TP + TN}{TP + TN + FP + FN} = \frac{290 + 292}{651} = 89.40\% \\
\text{Precision} &= \frac{TP}{TP + FP} = \frac{290}{290 + 46} = 86.31\% \\
\text{Recall} &= \frac{TP}{TP + FN} = \frac{290}{290 + 23} = 92.65\% \\
\text{Specificity} &= \frac{TN}{TN + FP} = \frac{292}{292 + 46} = 86.39\%
\end{align}

\textbf{Nhận xét:} Confusion Matrix cho thấy mô hình có khả năng phát hiện té ngã tốt với Recall đạt 92.65\%. Tuy nhiên, số lượng False Positive (46) và False Negative (23) cho thấy vẫn còn một số trường hợp khó phân loại:

\begin{itemize}
    \item \textbf{False Positive Analysis}: Các hoạt động sinh hoạt có tốc độ thay đổi tư thế nhanh (như ngồi xuống đột ngột, cúi người lấy đồ) có thể bị nhầm với té ngã do có đặc trưng động học tương tự.
    \item \textbf{False Negative Analysis}: Các cú té ngã chậm (slow descent fall) hoặc té ngã từ từ có đặc trưng không khác biệt nhiều so với ADL thông thường, dẫn đến bị bỏ sót.
\end{itemize}

\vskip 0.5cm

\begin{figure}[hbtp]
    \centering
    %\includegraphics[width=0.7\textwidth]{figures/chapter3/roc_curve_5fold_cv.png}
    \caption[ROC Curve 5-Fold CV]{Đường cong ROC (Receiver Operating Characteristic) từ 5-Fold Cross-Validation}
    \label{fig:roc_curve}
\end{figure}

Hình \ref{fig:roc_curve} thể hiện đường cong ROC với AUC (Area Under Curve) = 0.9569. Đây là chỉ số quan trọng đánh giá khả năng phân biệt tổng thể của mô hình.

\textbf{Nhận xét ROC Curve:}
\begin{itemize}
    \item \textbf{AUC = 0.9569}: Giá trị này cho thấy mô hình có khả năng phân biệt rất tốt giữa lớp Fall và No Fall. AUC gần 1.0 được coi là xuất sắc.
    \item \textbf{Đường cong ROC}: Đường cong ROC nằm xa đường chéo (random classifier) cho thấy mô hình hoạt động tốt hơn so với phân loại ngẫu nhiên ở hầu hết các ngưỡng phân loại.
    \item \textbf{Trade-off giữa TPR và FPR}: Đường cong cho thấy khi giảm ngưỡng phân loại để tăng TPR (phát hiện nhiều Fall hơn), FPR (báo động sai) cũng tăng. Ngưỡng tối ưu nên được chọn dựa trên yêu cầu cụ thể của ứng dụng.
\end{itemize}

Trong ứng dụng phát hiện té ngã, việc \textbf{minimize False Negative} (bỏ sót té ngã) thường quan trọng hơn việc giảm False Positive, vì bỏ sót một cú té ngã có thể dẫn đến hậu quả nghiêm trọng. Do đó, ta có thể chọn ngưỡng phân loại thấp hơn để đạt Recall cao hơn, chấp nhận một mức False Positive nhất định.

\vskip 0.5cm

\begin{figure}[hbtp]
    \centering
    %\includegraphics[width=0.7\textwidth]{figures/chapter3/pr_curve_5fold_cv.png}
    \caption[PR Curve 5-Fold CV]{Đường cong Precision-Recall từ 5-Fold Cross-Validation}
    \label{fig:pr_curve}
\end{figure}

Hình \ref{fig:pr_curve} thể hiện đường cong Precision-Recall với AP (Average Precision) = 0.9536. PR Curve đặc biệt hữu ích cho bài toán có sự mất cân bằng lớp.

\textbf{Nhận xét PR Curve:}
\begin{itemize}
    \item \textbf{AP = 0.9536}: Giá trị này cho thấy mô hình duy trì được Precision cao ngay cả khi Recall tăng, đây là dấu hiệu tốt.
    \item \textbf{So sánh với Baseline}: Đường baseline (nét đứt màu đỏ) thể hiện tỷ lệ Fall trong tập dữ liệu. Đường cong PR của mô hình nằm cao hơn baseline ở hầu hết các điểm, cho thấy mô hình có giá trị thực sự.
    \item \textbf{Precision cao ở Recall thấp}: Khi Recall thấp (chỉ phát hiện các mẫu dễ nhất), Precision đạt gần 100\%. Điều này cho thấy các mẫu "dễ" (có đặc trưng Fall rõ ràng) được phân loại rất chính xác.
    \item \textbf{Độ suy giảm Precision}: Khi tăng Recall để phát hiện thêm các trường hợp khó, Precision giảm dần. Tuy nhiên, mức độ suy giảm vẫn ở mức chấp nhận được.
\end{itemize}

Đối với bài toán phát hiện té ngã, PR Curve cho thấy mô hình có thể đạt được:
\begin{itemize}
    \item \textbf{Recall cao (>90\%)} với Precision còn khoảng 85-90\%
    \item \textbf{Precision cao (>95\%)} với Recall còn khoảng 80-85\%
\end{itemize}

Tùy vào yêu cầu ứng dụng (ưu tiên phát hiện đủ hay ưu tiên độ chính xác), ta có thể điều chỉnh ngưỡng phân loại phù hợp.

\vskip 0.5cm

\textbf{Tổng hợp nhận xét về ba đồ thị đánh giá:}

Ba đồ thị -- Confusion Matrix, ROC Curve và PR Curve -- cùng xác nhận rằng hệ thống HybridFallTransformer đạt hiệu suất tốt trên tập dữ liệu kiểm tra:

\begin{enumerate}
    \item \textbf{Độ chính xác tổng thể}: 89.40\% accuracy từ Confusion Matrix tổng hợp, cao hơn mức trung bình 81.26\% từ 5-Fold CV, cho thấy mô hình hoạt động ổn định.
    
    \item \textbf{Khả năng phân biệt xuất sắc}: AUC = 0.9569 từ ROC Curve cho thấy mô hình có khả năng phân biệt xuất sắc giữa Fall và No Fall.
    
    \item \textbf{Cân bằng Precision-Recall tốt}: AP = 0.9536 từ PR Curve xác nhận mô hình duy trì được độ chính xác cao khi tăng độ phủ.
    
    \item \textbf{Recall cao (92.65\%)}: Từ Confusion Matrix, Recall cao cho thấy mô hình ít bỏ sót các trường hợp té ngã thực sự -- điều rất quan trọng trong ứng dụng thực tế.
\end{enumerate}

Tuy nhiên, vẫn còn một số điểm cần cải thiện:
\begin{itemize}
    \item \textbf{46 False Positive}: Báo động sai có thể gây phiền toái cho người dùng nếu xảy ra thường xuyên.
    \item \textbf{23 False Negative}: Bỏ sót 23 cú té ngã vẫn là con số đáng lo ngại trong môi trường y tế.
\end{itemize}

\begin{itemize}
    \item \textbf{Fold 1 và Fold 2} hoạt động tốt nhất với Acc > 83\%
    \item \textbf{Fold 3} có kết quả thấp nhất, có thể do sự phân bố khó của dữ liệu trong fold này
    \item \textbf{Fold 4 và Fold 5} ở mức trung bình với Acc ~81\%
\end{itemize}

Các nguyên nhân chính gây ra lỗi:

\begin{itemize}
    \item \textbf{False Positive}: Các hoạt động sinh hoạt có tốc độ thay đổi tư thế nhanh (như ngồi xuống đột ngột) có thể bị nhầm với té ngã.
    \item \textbf{False Negative}: Các cú té ngã chậm (slow descent fall) có đặc trưng động học không khác biệt nhiều so với ADL thông thường.
\end{itemize}

\subsection{So sánh với các phương pháp khác}
\label{sec:comparison}

Bảng \ref{tab:method_comparison} so sánh HybridFallTransformer với các phương pháp tiêu biểu trên cùng bộ dữ liệu:

\begin{table}[h]
\centering
\caption{So sánh với các phương pháp khác}
\label{tab:method_comparison}
\begin{tabular}{|l|c|c|c|c|c|}
\hline
\textbf{Phương pháp} & \textbf{Data} & \textbf{Acc (\%)} & \textbf{F1 (\%)} & \textbf{AUC} & \textbf{Params} \\
\hline
Threshold-based \cite{} & Sensor & 75-85 & 70-80 & 0.78-0.85 & < 1K \\
\hline
3D-CNN \cite{} & RGB & 85-92 & 83-90 & 0.90-0.95 & 8-12M \\
\hline
LSTM \cite{} & Skeleton & 78-88 & 75-85 & 0.82-0.90 & 1-3M \\
\hline
GCN \cite{khani2024skeleton} & Skeleton & 88-92 & 86-90 & 0.92-0.95 & 2-5M \\
\hline
\textbf{HybridFallTransformer} & \textbf{Skeleton} & \textbf{81.26} & \textbf{80.05} & \textbf{0.847} & \textbf{1.8M} \\
\hline
\end{tabular}
\end{table}

HybridFallTransformer đạt kết quả ở mức trung bình-cao so với các phương pháp hiện có. Điểm mạnh của phương pháp này là kiến trúc Transformer với khả năng nắm bắt phụ thuộc dài hạn trong chuỗi khung xương. Tuy nhiên, kết quả còn có thể được cải thiện với:

\begin{itemize}
    \item Tăng kích thước bộ dữ liệu huấn luyện
    \item Tối ưu siêu tham số (hyperparameter tuning)
    \item Áp dụng data augmentation cho dữ liệu skeleton
\end{itemize}

%-------------------------------------------------

\section{Đánh giá khả năng thời gian thực}
\label{sec:realtime_evaluation}

\subsection{FPS và độ trễ suy luận}
\label{sec:fps_latency}

Bảng \ref{tab:fps_latency} trình bày hiệu suất suy luận của hệ thống trên các nền tảng phần cứng:

\begin{table}[h]
\centering
\caption{Hiệu suất suy luận trên các nền tảng phần cứng}
\label{tab:fps_latency}
\begin{tabular}{|l|c|c|c|}
\hline
\textbf{Hardware} & \textbf{FPS (batch=1)} & \textbf{Latency (ms)} & \textbf{Batch FPS} \\
\hline
Kaggle GPU (P100) & ~500 & ~2.0 & ~1600 \\
\hline
CPU (Modern) & 50-100 & 10-20 & 200-400 \\
\hline
Raspberry Pi 4 & 8-12 & 80-120 & 30-50 \\
\hline
\end{tabular}
\end{table}

Trên GPU, hệ thống đạt tốc độ rất cao với hơn 500 FPS cho single sample inference. Trên Raspberry Pi 4, hệ thống đạt 8-12 FPS -- đủ để xử lý video từ camera giám sát thông thường (thường 15-30 FPS).

Độ trễ end-to-end từ khung hình được chụp đến khi có kết quả phân loại là khoảng 2ms trên GPU và 80-120ms trên Raspberry Pi 4. Với camera 15 FPS (mỗi khung cách nhau 66.7ms), hệ thống có thể xử lý mọi khung hình trên thiết bị biên.

\subsection{Triển khai trên thiết bị biên Edge AI}
\label{sec:edge_deployment}

Đánh giá chi tiết trên Raspberry Pi 4 Model B 4GB:

\begin{itemize}
    \item \textbf{FPS đạt được}: 8-12 FPS (xử lý video từ camera 15 FPS)
    \item \textbf{Bộ nhớ sử dụng}: phụ thuộc vào kích thước batch và mô hình
    \item \textbf{Tổng kích thước mô hình}: ~7-10MB (YOLO + Transformer)
\end{itemize}

Với kết quả này, hệ thống có thể triển khai trên Raspberry Pi 4 để giám sát thời gian thực trong môi trường gia đình hoặc phòng bệnh nhân với độ trễ chấp nhận được.

\begin{figure}[h]
    \centering
    %\includegraphics[width=0.7\textwidth]{figures/chapter3/edge_benchmark.png}
    \caption[So sánh FPS trên các thiết bị]{So sánh FPS trên các thiết bị biên}
    \label{fig:edge_benchmark}
\end{figure}

Mặc dù Raspberry Pi 4 không đạt được tốc độ khung hình thời gian thực hoàn chỉnh cho camera 30 FPS, hệ thống vẫn đáp ứng yêu cầu ứng dụng thực tế vì:

\begin{enumerate}
    \item Camera giám sát trong nhà thường có tốc độ 15-25 FPS
    \item Việc bỏ qua 2-3 khung hình không ảnh hưởng đáng kể đến khả năng phát hiện vì té ngã kéo dài nhiều khung hình
    \item Với 8-12 FPS, độ trễ trung bình giữa hai lần phân loại là dưới 125ms
\end{enumerate}

%-------------------------------------------------

\section{Kết luận chương}
\label{sec:chapter_summary}

Chương này đã trình bày chi tiết kiến trúc, thiết lập thực nghiệm và kết quả đánh giá của hệ thống HybridFallTransformer.

\subsection{Ưu điểm của hệ thống}

\begin{enumerate}
    \item \textbf{Hiệu năng phân loại ở mức chấp nhận được}: HybridFallTransformer đạt Accuracy trung bình 81.26\%, F1-score trung bình 80.05\% và AUC trung bình 0.8467 trên tập dữ liệu tổng hợp CaucaFall + MCFD với phương pháp 5-Fold Cross-Validation.
    
    \item \textbf{Kết quả tương đối ổn định qua các fold}: Accuracy dao động từ 77.69\% đến 83.21\%, trong khi F1-score dao động từ 76.80\% đến 82.81\%. Điều này cho thấy mô hình duy trì khả năng phân loại tương đối ổn định trên các tập kiểm tra khác nhau.
    
    \item \textbf{Khai thác dữ liệu khung xương thay vì ảnh RGB trực tiếp}: Hệ thống sử dụng đặc trưng pose gồm chuỗi vector khung xương 60 chiều, giúp giảm phụ thuộc vào thông tin nền ảnh và hạn chế xử lý trực tiếp dữ liệu hình ảnh nhạy cảm.
    
    \item \textbf{Kiến trúc mô hình tương đối gọn}: Khối HybridFallTransformer có khoảng 1.8M tham số, thấp hơn so với nhiều mô hình xử lý trực tiếp chuỗi ảnh RGB. Tuy nhiên, con số này chỉ tính riêng mô hình phân loại chuỗi thời gian, không bao gồm YOLOv11-Pose ở bước trích xuất đặc trưng.
    
    \item \textbf{Khả năng triển khai trên thiết bị biên}: Hệ thống đạt 8-12 FPS trên Raspberry Pi 4, đủ để xử lý video từ camera giám sát thông thường trong môi trường trong nhà.
    
    \item \textbf{Tốc độ xử lý cao trên GPU}: Trên GPU đạt >500 FPS với độ trễ chỉ 2ms, đáp ứng yêu cầu xử lý thời gian thực.
\end{enumerate}

\subsection{Nhược điểm của hệ thống}

\begin{itemize}
    \item \textbf{Kết quả chưa đạt mức state-of-the-art}: Hiệu năng của HybridFallTransformer chưa vượt trội so với một số phương pháp mạnh như 3D-CNN (85-92\%) hoặc GCN (88-92\%).
    
    \item \textbf{Hiệu suất biến động giữa các folds}: Fold 3 đạt kết quả thấp hơn đáng kể so với các folds khác (Accuracy 77.69\% so với mức 81-83\% của các folds khác), cho thấy mô hình có thể nhạy cảm với phân bố dữ liệu trong từng fold.
    
    \item \textbf{Bộ dữ liệu huấn luyện còn hạn chế}: Với chỉ 651 mẫu, tập dữ liệu huấn luyện tương đối nhỏ so với các nghiên cứu khác, có thể ảnh hưởng đến khả năng tổng quát hóa của mô hình.
    
    \item \textbf{Phụ thuộc vào chất lượng trích xuất khung xương}: Hiệu suất của hệ thống phụ thuộc vào độ chính xác của YOLOv11-Pose trong việc trích xuất tọa độ các điểm khớp. Trong điều kiện ánh sáng yếu, occlusion hoặc nhiều người trong khung hình, chất lượng trích xuất có thể giảm.
    
    \item \textbf{Không xử lý được các trường hợp té ngã từ từ (slow descent)}: Các cú té ngã có tốc độ thay đổi chậm có đặc trưng động học không khác biệt nhiều so với hoạt động sinh hoạt thông thường, dẫn đến nguy cơ bị bỏ sót.
    
    \item \textbf{Số lượng False Positive còn đáng kể}: Với 46 trường hợp báo động sai trên 651 mẫu, hệ thống có thể gây phiền toái cho người dùng nếu triển khai trong thực tế.
\end{itemize}

%-------------------------------------------------
\section{Hướng phát triển}
\label{sec:future_work}

Dựa trên các hạn chế đã nhận diện, đề tài đề xuất ba hướng phát triển chính.

\subsection{Tối ưu mô hình và Edge AI}
\label{sec:future_edge}

\begin{enumerate}
    \item \textbf{Quantization và Pruning}: áp dụng các kỹ thuật nén mô hình như INT8 quantization và structured pruning để giảm kích thước và tăng tốc độ suy luận trên thiết bị biên.
    
    \item \textbf{Tối ưu cho Edge}: sử dụng TensorRT hoặc ONNX Runtime để tối ưu hóa mô hình trên thiết bị Edge, hướng tới đạt $30+$ FPS trên Raspberry Pi.
    
    \item \textbf{Edge-Cloud Hybrid}: thiết kế kiến trúc hybrid cho phép xử lý trên Edge khi điều kiện mạng không thuận lợi, đồng thời tự động chuyển lên Cloud khi cần tăng cường độ chính xác.
    
    \item \textbf{Kết hợp với TinyML}: nghiên cứu khả năng triển khai mô hình nhẹ hơn nữa trên vi điều khiển IoT.
\end{enumerate}

\subsection{Mở rộng bài toán và dữ liệu}
\label{sec:future_data}

\begin{enumerate}
    \item \textbf{Bộ dữ liệu thực tế}: thu thập thêm dữ liệu từ môi trường bệnh viện, viện dưỡng lão và hộ gia đình với sự đồng ý của người tham gia. Dữ liệu nên tập trung vào các tình huống khó như ánh sáng yếu, occlusion một phần và nhiều người trong khung hình.
    
    \item \textbf{Mở rộng phân loại đa lớp}: thay vì phân loại nhị phân Fall/ADL, mở rộng thành phân loại đa lớp với các nhãn như Standing, Walking, Sitting, Lying, Bending, Falling Forward, Falling Backward và Slipping, nhằm cung cấp thông tin chi tiết hơn về tình huống cho đội ngũ y tế.
    
    \item \textbf{Xử lý multi-person}: nghiên cứu kiến trúc xử lý đồng thời nhiều người trong khung hình, phù hợp với môi trường bệnh viện có nhiều bệnh nhân.
    
    % \item \textbf{Data Augmentation}: áp dụng các kỹ thuật augmentation cho dữ liệu skeleton để tăng tính đa dạng và độ ổn định của mô hình.
\end{enumerate}

%-------------------------------------------------

Nhìn chung, kết quả thực nghiệm cho thấy HybridFallTransformer có khả năng nhận dạng té ngã ở mức khá trên tập dữ liệu đã xây dựng. Mặc dù mô hình chưa vượt trội tuyệt đối so với một số phương pháp mạnh như 3D-CNN hoặc ST-GCN, hướng tiếp cận dựa trên khung xương vẫn có ưu điểm về tính gọn nhẹ, khả năng bảo vệ quyền riêng tư và phù hợp với bài toán theo dõi té ngã trong môi trường trong nhà.

Các kết quả này xác nhận tính khả thi của việc sử dụng kiến trúc Hybrid Transformer cho bài toán phát hiện té ngã trên thiết bị biên, mở ra hướng triển khai thực tế trong các cơ sở chăm sóc sức khỏe người cao tuổi.

%=================================================
